\begin{table}[H]
\centering
\caption{Nama Daerah Perbandingan Cluster Kompromi STATIS vs IPM (UNDP) (2024)} \label{tab:app2comp_ipm_name_2024}
\begin{adjustbox}{max width=\textwidth}
\begin{tabular}{|p{4cm}|p{10cm}|}
\hline
\textbf{Kategori} & \textbf{Daftar Wilayah} \\
\hline
Kompromi 1 (Rendah) \& Sangat Tinggi & - \\
\hline
Kompromi 1 (Rendah) \& Tinggi & Kep. Seribu, Ciamis, Kota Banjar \\
\hline
Kompromi 1 (Rendah) \& Sedang & Cianjur, Garut, Tasikmalaya, Majalengka, Indramayu, Pangandaran, Pandeglang, Lebak, Serang \\
\hline
Kompromi 1 (Rendah) \& Rendah & - \\
\hline
Kompromi 2 (Sedang) \& Sangat Tinggi & - \\
\hline
Kompromi 2 (Sedang) \& Tinggi & Bogor, Bandung, Kuningan, Cirebon, Sumedang, Subang, Purwakarta, Karawang, Bandung Barat, Kota Sukabumi, Kota Cirebon, Kota Tasikmalaya, Tangerang, Kota Cilegon, Kota Serang \\
\hline
Kompromi 2 (Sedang) \& Sedang & Sukabumi \\
\hline
Kompromi 2 (Sedang) \& Rendah & - \\
\hline
Kompromi 3 (Tinggi) \& Sangat Tinggi & Kota Jakarta Selatan, Kota Jakarta Timur, Kota Jakarta Pusat, Kota Jakarta Barat, Kota Jakarta Utara, Kota Bandung, Kota Bekasi, Kota Depok, Kota Cimahi, Kota Tangerang, Kota Tangerang Selatan \\
\hline
Kompromi 3 (Tinggi) \& Tinggi & Bekasi, Kota Bogor \\
\hline
Kompromi 3 (Tinggi) \& Sedang & - \\
\hline
Kompromi 3 (Tinggi) \& Rendah & - \\
\hline
\end{tabular}
\end{adjustbox}
\end{table}